% Chapter 6 - Conclusion and Future Work

\section{Conclusion and Future Work}\label{sec:conclusion}

This thesis set out from a single question, posed in \autoref{sec:introduction}:
whether a hand-driven, interactive 3D application could be built from nothing more
than a webcam and a commodity \acrshort{cpu}, with no headset, no depth camera, and
no instrumented gloves. The preceding chapters answer it in the affirmative and in
full, tracing a complete path from a single colour frame to four running interactive
scenes. What follows draws that work together, sets out honestly the limits of what
was achieved, and identifies the directions in which it could be carried further.

\subsection{Summary of Contributions}\label{subsec:conclusion_summary}

The central difficulty, identified in \autoref{subsec:gap}, was the gap between the
raw output of a model-free landmark detector and the canonical pose representation an
interactive application requires. MediaPipe \cite{zhang2020mediapipehands} recovers
21 landmarks per hand at interactive rate on commodity hardware, but those landmarks
are a flat list of points whose scale varies with camera distance, carry only
wrist-relative depth, and encode neither a joint hierarchy nor anatomically valid
ranges of motion~(\autoref{subsubsec:model-free}). The \acrshort{mano} model
\cite{MANO:SIGGRAPHASIA:2017} supplies exactly the missing structure --- a fixed
kinematic tree of joint rotations, independent of hand size and camera
distance~(\autoref{subsec:mano}) --- and bridging the two was the work of the thesis.

The bridge is the analytical AIK solver of \autoref{subsec:aik}, examined against the
four families of inverse-kinematics method surveyed in \autoref{subsec:sota_ik}.
Because the 21-landmark case is fully observed, the pose can be recovered in closed
form: a single \acrshort{svd}-based point-set alignment for the global orientation,
followed by one swing-only traversal of the kinematic tree~(\autoref{subsec:ik_impl}).
The result, as \autoref{subsubsec:sota_ik_summary} argues, is the only surveyed
approach that is simultaneously data-free, deterministic in per-frame cost, and
runnable on \acrshort{cpu} --- the qualities the deployment setting demands.

Around that solver, \autoref{sec:pipeline} described the Python pipeline that turns
each frame into a smoothed, scale-normalised pose: metric landmark extraction with
adaptive 1\euro{} smoothing~(\autoref{subsec:landmark_smoothing}), the \acrshort{mano}
forward model used for the rest-pose template and local
visualisation~(\autoref{subsec:forward_model}), and a \acrshort{udp} streaming layer
that serialises the solved rotations, a gesture label, and a placement anchor as a
\acrshort{json} datagram~(\autoref{subsec:streaming}). \autoref{sec:application} then
described the Unity application that consumes the stream: a receiver that reverses the
serialisation step for step~(\autoref{subsec:receiver}), a shared interaction layer
providing one gesture vocabulary, one body-anchored menu, and one pointing
component~(\autoref{subsec:interaction_layer}), and four scenes that put the same pose
to markedly different uses --- a companion-and-grappling bear, a vertex-by-vertex face
deformer, an editable city, and a painting tool whose drawings leave the application
as point clouds for the surface reconstruction of \autoref{subsec:nksr}.

That this pipeline is not merely functional but responsive is the quantitative
contribution of the work. The latency evaluation of \autoref{subsec:latency_eval}
measured a mean software latency of \SI{30.68}{\ms} over \num{20234}
frames, with a 99th percentile of \SI{40}{\ms} --- comfortably below the \SI{60}{\ms}
perceptibility threshold identified in \autoref{subsec:one_euro_filter}, and achieved
on a laptop \acrshort{cpu} with no \acrshort{gpu} involvement at any stage. The
recurring observation across all four scenes, framed throughout against the state of
the art~(\autoref{subsec:sota_interaction}), is therefore that each interaction is of
a kind that comparable systems deliver only with a dedicated peripheral, a
head-mounted display, or a controller, and each is achieved here from a single webcam.

\subsection{Limitations}\label{subsec:conclusion_limitations}

The system's strengths follow directly from a set of deliberate trade-offs, and those
trade-offs are equally the source of its limitations. The most consequential is one of
evaluation rather than design: the recovered pose was validated for latency but not
for accuracy. No benchmark figure quantifies how closely the streamed \acrshort{mano}
rotations reproduce ground-truth hand pose, and the system's correctness rests instead
on the geometric guarantees of the closed-form solver and on visual inspection of the
rendered avatars. This is the natural complement to the latency result and its most
obvious omission.

A second limitation is structural to the solver. AIK sets the twist component of every
finger joint to zero~(\autoref{subsec:aik}), a simplification that is physiologically
well grounded but that discards exactly the degree of freedom a learned method such as
HybrIK-X \cite{li2023hybrikx} is able to recover from image
features~(\autoref{subsubsec:sota_ik_neural}). For most finger articulation this loss
is negligible, but for poses in which axial rotation carries genuine information ---
thumb opposition above all --- the closed-form path cannot represent what it does not
solve for. As \autoref{subsubsec:sota_ik_summary} put it, HybrIK-X defines the accuracy
ceiling this system does not reach, while AIK defines the latency and hardware floor
within which it operates.

The reliance on a single RGB camera imposes its own limits. Depth is not measured but
inferred: the placement anchor combines the wrist's image-space position with a
bone-length-ratio depth cue~(\autoref{subsec:landmark_smoothing}), an estimate adequate
for plausible placement but not for metric reconstruction, and far short of the spatial
precision and tracking volume that the dedicated rigs of \autoref{subsubsec:sota_creative}
provide. A single viewpoint also offers no recourse when a hand leaves the frame or
turns away from the camera, and self-occlusion of the fingers degrades the landmarks
before the solver ever sees them. The grace windows of \autoref{subsec:interaction_layer}
absorb brief detector dropouts, but they are a tolerance, not a recovery mechanism.

Finally, the interaction layer is bounded by the detector it builds on. The application
draws on five discrete gestures from MediaPipe's recogniser~(\autoref{subsec:interaction_layer})
and tracks at most two hands; richer or user-defined gestures, and continuous rather
than discrete control signals, lie outside the present vocabulary.

\subsection{Directions for Future Work}\label{subsec:conclusion_future}

Each of these limitations suggests a corresponding line of further work, and they are
set out here in the same pipeline order the thesis has followed throughout.

The most immediate is the quantitative evaluation the present work lacks. Benchmarking
the recovered pose against established datasets such as FreiHAND or HO-3D, on which the
regressors of \autoref{subsubsec:model-based} report their accuracy, would place this
system on the same footing as the methods it is compared against and quantify the cost
of the zero-twist simplification directly. A controlled user study of the four
interactions would complement that geometric measure with evidence on the experience
itself.

At the estimation layer, the clearest opportunity is to narrow the accuracy gap without
surrendering the deployment target. A lightweight learned twist head, in the spirit of
HybrIK \cite{li2022hybrik}, could augment the closed-form swing with the one component
geometry alone cannot determine, running only when a \acrshort{gpu} is available and
falling back to pure AIK otherwise --- preserving the \acrshort{cpu}-only floor while
reaching toward the learned ceiling. Fusing a second camera or an optional depth sensor
would in turn restore the metric scale the monocular anchor can only approximate, for
users willing to add that hardware.

The model itself admits two refinements. Per-user shape estimation, which the present
pipeline fixes rather than solves, would fit the \acrshort{mano} template to the
individual hand and improve the proportions the solver normalises
against~(\autoref{subsec:mesh_recon}). More visibly, the Unity avatar is driven as a
skinned armature and so renders under plain linear blend skinning, without the
corrective pose blend shapes that \acrshort{mano} learns in order to suppress joint
collapse~(\autoref{subsec:mano}). Carrying the full \acrshort{mano} mesh and its pose
blend-shape term~(\autoref{eq:mano_bp}) through to the renderer would bring the finger
deformation of the displayed hand up to the fidelity the model is capable of.

Robustness is the natural next concern. An occlusion-aware front end along the lines of
HandOccNet \cite{park2022handoccnet}, or a temporal model that interpolates across
dropped frames rather than merely tolerating them, would extend reliable tracking into
the partial-occlusion and out-of-frame conditions a single camera currently handles
poorly.

At the interaction layer, a custom or dynamic gesture vocabulary would lift the
five-gesture limit, and the surface-reconstruction step that the painting tool feeds
offline~(\autoref{subsec:scene_paint}) could be brought inside the application, so that
a point cloud painted in front of the webcam is reconstructed into a mesh live by
NKSR \cite{huang2023nksr} rather than exported for separate processing. Beyond the
single machine, the loopback transport of \autoref{subsec:streaming} generalises
naturally to a networked stream, opening remote or multi-user interaction over the same
fire-and-forget channel.

\subsection{Concluding Remarks}\label{subsec:conclusion_closing}

The thread running through both the design and its limitations is a single trade, made
deliberately and consistently: accuracy and tracking fidelity are given up in exchange
for accessibility. The most capable systems surveyed in this thesis reach their
precision by adding hardware --- a depth camera, a head-mounted display, an optical
peripheral, a pair of controllers --- and in doing so raise the cost of entry to the
very interactions they enable. The system built here takes the opposite position. It
accepts the ceiling that a single RGB camera and a closed-form solver impose, and in
return asks the user for nothing they do not already own.

What the four scenes demonstrate, taken together, is that this floor is high enough to
be useful. A character flown by the hands, a face pulled out of shape, a city
rearranged in place, and a drawing lifted off the page into a reconstructed solid are
not coarse cursor tasks but fine-grained, three-dimensional, two-handed interactions
--- and they run, responsively, on the camera already built into a laptop. The
accessibility this affords was the motivation with which the thesis
opened~(\autoref{sec:introduction}); that it can be delivered at interactive latency,
with no hardware beyond the commodity, is its conclusion.
